\section{Introduction}
\label{sec:introduction}

Queue-based execution decouples work production from work execution. A producer
submits a task instance to a queue, and one or more workers later reserve,
execute, retry, acknowledge, or fail that work. Although this pattern is often
presented as a background-job abstraction, production queues frequently execute
operationally significant work: replication, repair, cache rebuilds, imports,
media processing, external-service calls, validation, notification delivery, and
other asynchronous operations. In such systems, a worker is an execution boundary
rather than merely a message consumer.

This paper studies profiling at the level of logical task attempts inside
queue-worker systems. This level is narrower than a worker process, container,
node, or queue, but broader than an individual function call or trace span. It is
the level at which a particular execution of a queued task consumes resources,
holds bounded execution capacity, produces outcomes, and can be retried or
failed. Coarser observations are useful, but they do not by themselves explain
which task kind or attempt produced the observed behavior.

The central problem is that task behavior is conditional. A task kind and a
payload-size bucket may provide a useful baseline, but they do not fully
determine resource demand or worker occupancy. Two attempts of the same
replication task kind may process similar payloads and still differ because one
runs against healthy peers while another runs under degraded peer health,
unstable network paths, disk pressure, or retry backoff. The task kind has not
changed; the execution context has changed, and that context may amplify only
some cost dimensions.

This paper uses a role-specific view of profiling. Execution context is treated
as input-side state. Active resource demand and worker occupancy are
attempt-attributed cost-side estimates. Uncertainty and operational risk are
downstream interpretations of estimates, residuals, missing information, and
consequences. These roles are related, but they are not synonyms. The purpose of
the model is to preserve these distinctions long enough for downstream systems
to make explicit decisions.

Existing work motivates parts of this view. Cluster managers and fairness
policies reason about multiple resource dimensions and observed usage
\cite{verma2015borg,ghodsi2011drf}. Workload studies show heterogeneity,
dynamic behavior, and task-usage variation at scale
\cite{zhang2011taskusage,reiss2012heterogeneity,cheng2018alibaba}. Workload
classification systems infer behavior rather than relying only on static
resource assumptions \cite{delimitrou2013paragon,delimitrou2014quasar}. Large
distributed systems are often shaped by tail behavior rather than average case
alone \cite{dean2013tail}. Observability systems provide essential metrics,
traces, logs, and profiles, but they introduce overhead and require budgeted
sampling or instrumentation policies
\cite{sigelman2010dapper,opentelemetrySampling,opentelemetryPerformance}. These
lines of work do not remove the need for a queue-worker-specific model that
connects logical tasks, execution context, cost estimates, uncertainty, and
profiling policy.

This paper proposes \emph{Task Behavior Graphs} as that model. A Task Behavior
Graph is a task-kind- or task-family-level representation that connects
canonical features to resource-demand, worker-occupancy, uncertainty, and
risk-relevant outputs through explicit influence assumptions. The graph is not
treated as proof of causality. It is a structured modeling artifact: it exposes
which features are believed to matter, which cost dimensions they influence, and
where the model remains uncertain.

The paper makes five contributions. First, it frames queue-worker task profiling
as conditional behavior estimation rather than raw telemetry collection. Second,
it separates intrinsic task features, execution context, active resource demand,
worker occupancy, uncertainty, and operational risk into distinct model roles.
Third, it introduces Task Behavior Graphs as a representation for connecting
these roles at task-kind or task-family granularity. Fourth, it describes how
raw observations can be mapped into canonical features through descriptors,
manifests, and domain-specific extensions. Fifth, it outlines how profiling
policy can be selected from task behavior, uncertainty, consequence, and
instrumentation budget.

The remainder of the paper is organized as follows.
Section~\ref{sec:terminology-notation} defines the core terminology and
notation. Section~\ref{sec:background} provides the technical background and
motivating distinctions. Section~\ref{sec:problem-statement} states the
profiling problem. Sections~\ref{sec:task-cost-model}--\ref{sec:task-behavior-graphs}
introduce the cost model, influence domains, task taxonomy, and graph
representation. Sections~\ref{sec:metrics-signals-features}--\ref{sec:profiling-policy-selection}
discuss feature mapping, profiling methods, observation points, uncertainty,
risk, and policy selection. Section~\ref{sec:task-behavior-manifests} describes
manifests and domain packs. The remaining sections discuss case studies,
evaluation, limitations, related work, and conclusions.
